\documentclass[10pt]{article}
\usepackage[margin=20mm]{geometry}
\usepackage{amsmath,amssymb,fontspec,unicode-math,xurl,hyperref}
\IfFontExistsTF{DejaVu Serif}{\setmainfont{DejaVu Serif}}{\setmainfont{Latin Modern Roman}}
\setmathfont{latinmodern-math.otf}
\setlength{\parindent}{0pt}\setlength{\parskip}{6pt}\sloppy
\title{Scientific review of invariant LLM unlearning}
\author{}
\date{Initial review, 25 October 2026}
\begin{document}
\maketitle
Abstract: The reviewed ILU paper supplies useful empirical evidence of resistance to the evaluated downstream fine-tuning tasks. A task-vector sign inference is invalid, and several descriptions of the numerical results and stopping protocol need correction. These are supported local issues (Minor); the scalar invariance surrogate, incomplete released implementation, and broader robustness interpretation require additional evidence. No failure of the reported six-task performance advantage is established here. Scientific adjudication remains pending.\par
\section*{Identity, scope, and principal result}
This report concerns Invariance Makes LLM Unlearning Resilient Even to Unanticipated Downstream Fine-Tuning, by Changsheng Wang, Yihua Zhang, Jinghan Jia, Parikshit Ram, Dennis Wei, Yuguang Yao, Soumyadeep Pal, Nathalie Baracaldo, and Sijia Liu, ICML 2025, PMLR 267, pp. 65464–65479 [paper]. The assignment is 0228/I25-43597, OpenReview identifier `x2lm33kdrZ`. The reviewed object is the staged 16-page publisher PDF, retrieved on 25 October 2026. Page references below use its PDF pages, coinciding with its printed 1–16 numbering.\par
The complete main paper and appended material were read. Checks cover Eqs. (1)–(5), the task-vector analysis, the definitions of forget quality (FQ), robust accuracy (RA), and fine-tuning accuracy (FA), Tables 1–4/A1–A2, and the experimental qualifications. The preserved repository was inspected statically, including the registered training methods, objective, configuration, and evaluation dispatch. All 58 archived members had verified archive/checkout correspondence. No author code or model was executed; no empirical replication is claimed.\par
The central empirical comparison remains supported by the printed results within its tested scope. For example, Table 2 reports average RA of 0.65 for RMU+ILU and 0.42 for RMU. RA is explicitly a mean over three training milestones, not a worst-case guarantee. Figures 3/A2 separately give final-epoch FQ and Fig. 5 gives trajectories, so this definition alone does not invalidate the resilience evidence.\par
\section*{Necessary distinctions}
ILU adds a penalty to an unlearning objective: \[\ell_u(\theta)+\lambda\sum_i s_i(\theta)^2,\qquad
 s_i(\theta)=\left.\partial_w\ell_i(w\circ\theta)\right|_{w=1}.\] Here \(w\) is the paper’s virtual scalar predictor. The paper explicitly describes IRMv1 as a relaxation and its stationarity condition as necessary, not sufficient. A stationary scalar direction is distinct from stationarity in all model parameters, and either is distinct from preservation of forgetting under arbitrary future optimization.\par
In Section 5, let \[u=\theta_u-\theta_o,\quad f=\theta_{ft}-\theta_o,\quad
 d=\theta_{uft}-\theta_u.\] The vectors \(u\) and \(f\) start at the original model; \(d\) starts at the unlearned model. The final displacement from the original model is \(u+d\). Cosines require nonzero vectors. The paper’s sign argument assumes \(\cos(u,f)<0\), not that \(f\) is exactly a negative multiple of \(u\).\par
\section*{Supported problems}
Problem 1 (The two cosine-sign criteria are not opposites; Minor).  Section 5, p. 6, interprets \(\cos(d,u)\geq0\) as preserving the unlearning direction and, conversely, \(\cos(d,f)\geq0\) as misalignment with unlearning. The second implication does not follow from \(\cos(u,f)<0\).\par
Complete check. Choose \[u=(1,0),\qquad f=(-1/2,\sqrt3/2),\qquad
 d=u+f=(1/2,\sqrt3/2).\] All vectors have norm one and satisfy the stipulated \(\cos(u,f)=-1/2<0\). Nevertheless, \[\cos(d,u)=1/2,\qquad \cos(d,f)=1/2.\] Thus both alleged opposing diagnostics hold simultaneously. This is not an extraction artifact: the paper itself reports the same sign pattern for ILU, namely \(\cos(d,u)=0.09\) and \(\cos(d,f)=0.3554\), while \(\cos(u,f)=-0.64\). ◻\par
There is also a distinction between an update and the final displacement: \[u^\top(u+d)=\|u\|^2+\|u\|\|d\|\cos(d,u).\] A negative update cosine need not make the final displacement oppose \(u\). For example, \(d=-0.1u\) leaves \(u+d=0.9u\). With the reported NPO cosine \(-0.41\), a negative final projection would additionally require \(\|d\|/\|u\|>1/0.41\approx2.439\). This ratio is not supplied.\par
The correct local statement is that \(\cos(d,u)\geq0\) increases or preserves the projection of \(u+d\) onto \(u\). It does not establish a forgetting property of the resulting network. Positive alignment with \(f\) implies negative alignment with \(u\) only under stronger conditions, such as exact antiparallel reference vectors with strictly positive alignment. Those conditions are absent here.\par
The reported cosine triples are geometrically admissible: their Gram determinants are approximately 0.080604 (NPO) and 0.415049 (ILU). No fabricated measurement is inferred. The defect affects the explanatory sign diagnostic and its causal interpretation, rather than a theorem or optimization guarantee on which ILU’s objective depends. This is the reason for Minor calibration; unchanged empirical scores alone would not downgrade a defective core theorem. Evidence: p. 6, `s06.png`, and independent `check.py`/`calc.json`.\par
Problem 2 (Descriptions overstate uniform forgetting and utility preservation; Minor).  The p. 8 narrative says MUSE-News ILU retains zero VerbMem across all fine-tuning settings, whereas Table A1, p. 15, reports 0.46 after GSM8K and 2.07 after QQP. The p. 7 claim that FA improves relative to either baseline is also not uniform in its own comparisons.\par
Table 2 gives average FA 82.43 for RMU and 82.32 for RMU+ILU(GSM8K), a decrease of 0.11 percentage point. This small difference, without uncertainty estimates, does not establish a significant degradation. The larger MUSE trade-offs should nevertheless be explicit: before fine-tuning, retained-set KnowMem on MUSE-Books falls from 57.19 for NPO to 45.20 for ILU, a difference of \(-11.99\) points; it remains lower in all six post-fine-tuning settings. MUSE-News WinoGrande FA is 67.40 versus 59.00, a difference of \(-8.40\) points.\par
Counterevidence bounds this issue. Appendix B.3 already says zero in nearly all settings; the nonzero ILU VerbMem scores remain far below NPO’s 35.38 and 33.46. Other utility comparisons remain favorable or close, so general model collapse has not been established. The defensible description is improved measured robustness with dataset-dependent utility trade-offs and small residual memorization. Utility-matched experiments would be needed to separate robustness gains from capability loss in the affected MUSE comparisons. Evidence: Tables 2/A1 and `calc.json`; the exact tabulated values, not a rerun, support this conclusion.\par
Problem 3 (The convergence rule is inconsistent; Minor).  Section 6.1, p. 7, specifies less than 1\% FA change over three consecutive epochs; Appendix A.2, p. 13, specifies two consecutive epochs and applies that criterion to all six tasks.\par
These are distinct stopping rules. For a run with two sufficiently small successive changes followed by a larger change, one rule can stop while the other continues. Since RA averages FQ at the first quartile, median, and final epoch, the stopping horizon can change all three checkpoints. This is a reproducibility defect in the stated protocol. The direction or size of any effect on the published RA values is not established. A single documented rule, actual per-run horizons, and checkpoint logs are required; choosing either printed rule without those records would not verify a repair. Evidence: `s07.png`, `s13.png`, and the bound page text.\par
\section*{Scope clarifications and unresolved evidence}
\subsection*{Scalar stationarity}
For the benign binary cross-entropy model \[\ell(w\theta)=\log(1+\exp(-w\theta)),\] at \(\theta=0,w=1\) one has \(\partial_w\ell=0\), hence zero regularization penalty, but \(\partial_\theta\ell=-1/2\). A gradient-descent step of size 0.1 moves \(\theta\) to 0.05. Thus a zero scalar penalty cannot certify full-parameter stationarity. This analytical example addresses an unrestricted interpretation of the mechanism, not the full ILU training objective or its measured performance. In light of the paper’s explicit relaxation language and absence of a universal theorem, stronger guarantees are Needs Evidence, not a manufactured failed theorem.\par
\subsection*{Released implementation and historical binding}
The preserved release [repo] states that code is ongoing. The method registry and command-line entrypoint expose NPO and other baseline methods, but no ILU or RMU implementation. The NPO loss in `src/unlearn/GA.py`, lines 114–129, consists of the forget loss plus the retain loss, without the invariance term. A repository-wide textual scan for the relevant invariance terms found hits only in the README; this negative search is corroborated by direct inspection of the registered objective path.\par
The WMDP configuration specifies NPO+FT, a mask, 500 steps, and WMDPALL for both datasets, whereas Appendix A.1 describes 2000 NPO steps, a biology forget set, and WikiText retain data. This establishes that this current configuration is not an adequate recipe for the paper’s described experiment. It does not establish what generated the historical results.\par
The actual evaluator calls the external `lm\_eval` wrapper for WMDP and MMLU. A separate TOFU forget-quality script and a standalone WMDP scorer are not proof of a metric error in that route. The evaluation harness revision, ILU implementation, downstream stopping/RA code, MUSE results, and timing logs are not bound by the inspected release. These remain Needs Evidence. Source excerpts and member SHA256/Git-blob bindings are in `code.json`; archives and their existing maps are preserved.\par
\subsection*{Numerical ambiguities and robustness limits}
Table 4, p. 9, prints original-model GSM8K FA as 0.75 alongside other FA percentages. The stated mean 75.02 equals the mean using 0.75. If that cell were intended to be 75, the mean would be 87.40. Both calculations are verified, but no run record establishes the intended value; no correction to 75 is asserted. Figure 6’s left/right baseline labels also merit checking against underlying records: the row near \(\lambda=0.1\) places FQ near 0.57 under the RMU label and 0.68 under NPO, unlike the main baseline levels. Different runs are possible, so a label swap is not established.\par
Table 3 discloses remaining targeted-relearning vulnerability: ILU FQ falls from 0.56 to 0.50 for NPO and from 0.68 to 0.54 for RMU. The stronger SAM comparison and future-work qualification are acknowledged by the authors. Arbitrary-task, adaptive, longer-horizon, or permanent-erasure interpretations exceed the tested scope. Reported TAR/ILU timing is about 63-fold, but the common hardware and execution binding cannot be independently verified from the available artifacts.\par
\section*{Dependence, corrections, and state}
The geometrical correction is exact but supplies no replacement theory of robust unlearning. The result descriptions can be corrected directly from Table A1. The stopping rule, historical implementation, ambiguous units, and matched-utility comparisons require records or new experiments. ArXiv metadata identifies a later v2 dated 10 October 2025 [arxiv]; its full content was not substituted for the frozen publisher PDF, and no claim is made that these observations persist in every later version or are previously unknown. The README’s arXiv link instead resolves to WAGLE, a different paper [wagle]; it was not treated as an ILU revision.\par
Item | Judgment | Repair verification | Adoption | Remaining work\par
Cosine inference | Minor | Projection identity proved | Unknown | Functional mechanism\par
Result descriptions | Minor | Table-based correction | Unknown | Matched utility\par
Stopping rule | Minor | No historical repair | Unknown | Horizons and logs\par
Scalar/release/units | Needs Evidence | None claimed | Unknown | Complete bindings\par
All scientific decisions remain pending; `human\_reviews=[]`. The report does not constitute acceptance.\par
\section*{Supporting checks}
and the preserved repository commit is\par
.\par
Independent `check.py` uses standard-library arithmetic and static text reads, produces `calc.json`/`code.json`, and renders original pages 4, 6, 7, 9, 13, and 15 with the existing Poppler executable. Those pages were visually checked against the extraction. The two-dimensional counterexample and scalar derivative are exact analytical checks; floating-point outputs merely confirm them.\par
No author code was executed.\par
There is no supplied accessible Space copy. The separate OpenReview supplement probe returned HTTP 403; this does not establish nonexistence. Empirical reproduction, full dependency-history verification, and exhaustive later-erratum searching remain unperformed.\par
9 C. Wang et al., Invariance Makes LLM Unlearning Resilient Even to Unanticipated Downstream Fine-Tuning, ICML 2025, PMLR 267, 65464–65479. Frozen publisher PDF. https://proceedings.mlr.press/v267/wang25en.html. OPTML-Group, Unlearn-ILU, preserved commit `0030928281846493d4f783db64aa45184b9dd814`, static inspection only. https://github.com/OPTML-Group/Unlearn-ILU. C. Wang et al., arXiv:2506.01339v2, 10 October 2025. Metadata consulted 25 October 2026; full later version not reviewed. https://arxiv.org/abs/2506.01339. J. Jia et al., WAGLE, arXiv:2410.17509v2, 12 April 2025. Metadata consulted only to resolve the release’s mismatched link. https://arxiv.org/abs/2410.17509.\par
\end{document}
